% GENERATED FILE. Edit canonical lessons, then run npm run generate:books.
% canonical-source-hash: fnv1a64:0124a8817ec78ad1
% canonical-lessons: AR-C207-faqat, AR-C207-jiddan, AR-C207-hatta, AR-C207-idha, AR-C207-kama

\chapter{Only, Very, Even, If, As}
\label{ch:ar-only-very-even-if-as}
\hlchaptermodality{207}

\begin{chapteropening}
\textbf{By the end of this chapter:} \emph{I can say only, very, even, if and as.}

Complete the last lesson of chapter 207: I can say only, very, even, if and as.
\end{chapteropening}

\section[\texorpdfstring{\ar{فقط} (faqaṭ)}{فقط (faqaṭ)}]{\ar{فقط} (faqaṭ) — only}

\label{lesson:AR-C207-faqat}



\begin{quote}
Before the new one: say the Arabic for joy, then the Arabic for fear.
\end{quote}

\subsection*{You'll want to know: \ar{فقط}}
\textbf{\ar{فقط}} — \emph{faqaṭ} — \textquotedblleft{}only\textquotedblright{}.

\begin{grammarlens}[title={What you've built}]
One more word for this chapter.
\end{grammarlens}

\subsection*{Your turn}
\begin{itemize}
\raggedright
  \item \emph{Say it:} \emph{faqaṭ}
  \item \emph{Say it:} \emph{faqaṭ}, once more
  \item \emph{Say it:} say \emph{faqaṭ}
  \item \emph{Recall:} say the Arabic for joy, then the Arabic for fear, then say \emph{faqaṭ} again
\end{itemize}

\begin{tcolorbox}[breakable,colback=teal!4,colframe=teal!35!black,title={Before you move on}]
What does \ar{فقط} mean? (\textquotedblleft{}Only\textquotedblright{}.) Say it once more.
\end{tcolorbox}
\section[\texorpdfstring{\ar{جدا} (jiddan)}{جدا (jiddan)}]{\ar{جدا} (jiddan) — very}

\label{lesson:AR-C207-jiddan}



\begin{quote}
Before the new one: say the Arabic for fear, then the Arabic for only.
\end{quote}

\subsection*{You'll want to know: \ar{جدا}}
\textbf{\ar{جدا}} — \emph{jiddan} — \textquotedblleft{}very\textquotedblright{}.

\begin{grammarlens}[title={What you've built}]
One more word for this chapter.
\end{grammarlens}

\subsection*{Your turn}
\begin{itemize}
\raggedright
  \item \emph{Say it:} \emph{jiddan}
  \item \emph{Say it:} \emph{jiddan}, once more
  \item \emph{Say it:} say \emph{jiddan}
  \item \emph{Recall:} say the Arabic for fear, then the Arabic for only, then say \emph{jiddan} again
\end{itemize}

\begin{tcolorbox}[breakable,colback=teal!4,colframe=teal!35!black,title={Before you move on}]
What does \ar{جدا} mean? (\textquotedblleft{}Very\textquotedblright{}.) Say it once more.
\end{tcolorbox}
\section[\texorpdfstring{\ar{حتى} (ḥattā)}{حتى (ḥattā)}]{\ar{حتى} (ḥattā) — even; until}

\label{lesson:AR-C207-hatta}



\begin{quote}
Before the new one: say the Arabic for only, then the Arabic for very.
\end{quote}

\subsection*{You'll want to know: \ar{حتى}}
\textbf{\ar{حتى}} — \emph{ḥattā} — \textquotedblleft{}even; until\textquotedblright{}.

\begin{grammarlens}[title={What you've built}]
One more word for this chapter.
\end{grammarlens}

\subsection*{Your turn}
\begin{itemize}
\raggedright
  \item \emph{Say it:} \emph{ḥattā}
  \item \emph{Say it:} \emph{ḥattā}, once more
  \item \emph{Say it:} say \emph{ḥattā}
  \item \emph{Recall:} say the Arabic for only, then the Arabic for very, then say \emph{ḥattā} again
\end{itemize}

\begin{tcolorbox}[breakable,colback=teal!4,colframe=teal!35!black,title={Before you move on}]
What does \ar{حتى} mean? (\textquotedblleft{}Even; until\textquotedblright{}.) Say it once more.
\end{tcolorbox}
\section[\texorpdfstring{\ar{إذا} (idhā)}{إذا (idhā)}]{\ar{إذا} (idhā) — if}

\label{lesson:AR-C207-idha}



\begin{quote}
Before the new one: say the Arabic for very, then the Arabic for even; until.
\end{quote}

\subsection*{You'll want to know: \ar{إذا}}
\textbf{\ar{إذا}} — \emph{idhā} — \textquotedblleft{}if\textquotedblright{}.

\begin{grammarlens}[title={What you've built}]
One more word for this chapter.
\end{grammarlens}

\subsection*{Your turn}
\begin{itemize}
\raggedright
  \item \emph{Say it:} \emph{idhā}
  \item \emph{Say it:} \emph{idhā}, once more
  \item \emph{Say it:} say \emph{idhā}
  \item \emph{Recall:} say the Arabic for very, then the Arabic for even; until, then say \emph{idhā} again
\end{itemize}

\begin{tcolorbox}[breakable,colback=teal!4,colframe=teal!35!black,title={Before you move on}]
What does \ar{إذا} mean? (\textquotedblleft{}If\textquotedblright{}.) Say it once more.
\end{tcolorbox}
\section[\texorpdfstring{\ar{كما} (kamā)}{كما (kamā)}]{\ar{كما} (kamā) — as, in the same way as}

\label{lesson:AR-C207-kama}



\begin{quote}
Before the new one: say the Arabic for even; until, then the Arabic for if.
\end{quote}

\subsection*{You'll want to know: \ar{كما}}
\textbf{\ar{كما}} — \emph{kamā} — \textquotedblleft{}as, in the same way as\textquotedblright{}.

\begin{grammarlens}[title={What you've built}]
One more word for this chapter.
\end{grammarlens}

\subsection*{Your turn}
\begin{itemize}
\raggedright
  \item \emph{Say it:} \emph{kamā}
  \item \emph{Say it:} \emph{kamā}, once more
  \item \emph{Say it:} say \emph{kamā}
  \item \emph{Recall:} say the Arabic for even; until, then the Arabic for if, then say \emph{kamā} again
\end{itemize}

\begin{tcolorbox}[breakable,colback=teal!4,colframe=teal!35!black,title={Before you move on}]
What does \ar{كما} mean? (\textquotedblleft{}As, in the same way as\textquotedblright{}.) Say it once more.
\end{tcolorbox}
